\subsection{莫雷尔-嘉当公式}

设 X 是 \(C^{r}\) 类流形（当 K 的特征 \textgreater0 时为有限维），L 是完备可赋范李代数。设 \(\alpha\) 是 X 上取值于 L 的 \(C^{r-1}\) 类 1 阶微分形式。设 \(x \in X\)。映射

\[
(u _ {1}, u _ {2}) \mapsto [ \alpha_ {x} (u _ {1}), \alpha_ {x} (u _ {2}) ]
\]

从 \(\mathrm{T}_x(\mathbf{X})\times \mathrm{T}_x(\mathbf{X})\) 到 L 的映射是 \(\mathrm{T}_x(\mathbf{X})\) 上取值于 L 的连续交替双线性形式。我们将其记为 \([\alpha ]_x^2\)，于是 \([\alpha ]^2\) 是 X 上取值于 L 的 2 阶微分形式。将 X 中 \(x\) 的一个开邻域识别为 Banach 空间的一个开子集，立即看出 \([\alpha ]^2\) 属于 \(\mathrm{C}^{r - 1}\) 类。若 \(\mathbf{X}'\) 是 \(\mathrm{C}^r\) 类流形且 \(f\colon \mathbf{X}'\to \mathbf{X}\) 是态射，则

\[
[ f ^ {*} (\alpha) ] ^ {2} = f ^ {*} ([ \alpha ] ^ {2}).\tag{28}
\]

设 \(\alpha, \beta\) 是 X 上取值于 L 的两个 \(C^{r-1}\) 类 1 阶微分形式。外积 \(\alpha \wedge \beta\)（《可微与解析流形》，R，7.8.2）是由 \(\alpha\) 与 \(\beta\) 构成的 X 上取值于 L 的 \(C^{r-1}\) 类 2 阶微分形式；有

\[
(\alpha \wedge \beta) _ {x} (u _ {1}, u _ {2}) = [ \alpha_ {x} (u _ {1}), \beta_ {x} (u _ {2}) ] - [ \alpha_ {x} (u _ {2}), \beta_ {x} (u _ {1}) ]\tag{29}
\]

对 \(u_{1}, u_{2}\) 属于 \(\mathrm{T}_x(\mathbf{X})\) 的情形，立即得到

\begin{enumerate}
\def\labelenumi{(\arabic{enumi})}
\setcounter{enumi}{29}
\tightlist
\item
\end{enumerate}

\[
[ \alpha + \beta ] ^ {2} = [ \alpha ] ^ {2} + [ \beta ] ^ {2} + \alpha \wedge \beta\tag{31}
\]

\[
\alpha \wedge \alpha = 2 [ \alpha ] ^ {2}.
\]

命题 51. 设 G 是李群（当 K 的特征 \textgreater0 时为有限维），\(a_{1},\ldots,a_{p}\) 是 L(G) 中的元素，F 是完备可赋范空间，\(\alpha\) 是 G 上取值于 F 的 \(p-1\) 阶微分形式。若 \(\alpha\) 左不变，则

\[
(d \alpha) _ {e} (a _ {1}, \dots , a _ {p}) = \sum_ {i <   j} (- 1) ^ {i + j} \alpha_ {e} ([ a _ {i}, a _ {j} ], a _ {1}, \dots , a _ {i - 1}, a _ {i + 1}, \dots , a _ {j - 1},
\]

\[
\left. a _ {j + 1}, \dots , a _ {p}\right).
\]

若 \(\alpha\) 右不变，则

\[
\begin{array}{l} (d \alpha) _ {e} (a _ {1}, \dots , a _ {p}) \\ = - \sum_ {i <   j} (- 1) ^ {i + j} \alpha_ {e} ([ a _ {i}, a _ {j} ], a _ {1}, \dots , a _ {i - 1}, a _ {i + 1}, \dots , a _ {j - 1}, a _ {j + 1}, \dots , a _ {p}). \end{array}
\]

假设 \(\alpha\) 左不变。根据《可微与解析流形》，R，8.5.7，有

\[
\begin{array}{l} (d \alpha) (L _ {a _ {1}}, \ldots , L _ {a _ {p}}) = \sum_ {i} (- 1) ^ {i - 1} L _ {a _ {i}} \alpha (L _ {a _ {1}}, \ldots , L _ {a _ {i - 1}}, L _ {a _ {i + 1}}, \ldots , L _ {a _ {p}}) \\ + \sum_ {i <   j} (- 1) ^ {i + j} \alpha ([ L _ {a _ {i}}, L _ {a _ {j}} ], L _ {a _ {1}}, \ldots , L _ {a _ {i - 1}}, L _ {a _ {i + 1}}, \ldots , L _ {a _ {j - 1}}, L _ {a _ {j + 1}}, \ldots , L _ {a _ {p}}). \end{array}
\]

但是 G 上的函数 \(\alpha (\mathrm{L}_{a_1},\dots ,\mathrm{L}_{a_{i - 1}},\mathrm{L}_{a_{i + 1}},\dots ,\mathrm{L}_{a_p})\) 左不变，因而为常数。所以

\[
\mathrm{L} _ {a _ {i}} \alpha (\mathrm{L} _ {a _ {1}}, \dots , \mathrm{L} _ {a _ {i - 1}}, \mathrm{L} _ {a _ {i + 1}}, \dots , \mathrm{L} _ {a _ {p}}) = 0.
\]

此外，\([L_{a_{i}}, L_{a_{j}}] = L_{\{a_{i}, a_{j}\}}\)（命题23），从而得到命题51的第一个公式。第二个公式可类似地建立，此时使用关系 \([R_{a_{i}}, R_{a_{j}}] = -R_{\{a_{i}, a_{j}\}}\)。

推论 1. 设 G 是李群（当 K 的特征 \textgreater0 时为有限维），且 ω、ω' 是 G 的左、右典范微分形式。则

\[
d \omega + [ \omega ] ^ {2} = 0, \quad d \omega^ {\prime} - [ \omega^ {\prime} ] ^ {2} = 0.
\]

根据命题51，

\[
\begin{array}{r l} (d \omega) _ {e} (a _ {1}, a _ {2}) & = - \omega_ {e} ([ a _ {1}, a _ {2} ]) = - [ a _ {1}, a _ {2} ] = - [ \omega_ {e} (a _ {1}), \omega_ {e} (a _ {2}) ] \\ & = - [ \omega ] _ {e} ^ {2} (a _ {1}, a _ {2}) \end{array}
\]

从而得到第一个公式。第二个公式可类似地建立。

推论 2. 假设 G 是有限维的。设 \((e_1, \ldots, e_n)\) 是 L(G) 的一组基，\((e_1^*, \ldots, e_n^*)\) 是对偶基，\((c_{ijk})\) 是相对于基 \((e_1, \ldots, e_n)\) 的 L(G) 的结构常数，并设 \(\omega_i\)（分别为 \(\omega_i^{\prime})\)）是 G 上取值于 K 的左（分别为右）不变微分形式，满足 \((\omega_i)_e = e_i^*\)（分别为 \((\omega_i^{\prime})_e = e_i^*)\)。则

\[
d \omega_ {k} + \sum_ {i <   j} c _ {i j k} \omega_ {i} \wedge \omega_ {j} = 0 (k = 1, 2, \dots , n)
\]

\[
d \omega_ {k} ^ {\prime} - \sum_ {i <   j} c _ {i j k} \omega_ {i} ^ {\prime} \wedge \omega_ {j} ^ {\prime} = 0 (k = 1, 2, \dots , n).
\]

若 \(r < s\)，

\[
\begin{array}{r l} \left(d \omega_ {k}\right) _ {e} (e _ {r}, e _ {s}) & = - \left(\omega_ {k}\right) _ {e} ([ e _ {r}, e _ {s} ]) \\ & = - \sum_ {i} c _ {r s i} (\omega_ {k}) _ {e} (e _ {i}) \\ & = - c _ {r s k} \\ & = - \sum_ {i <   j} c _ {i j k} (\omega_ {i} \wedge \omega_ {j}) _ {e} (e _ {r}, e _ {s}). \end{array}
\]

对于 \(\omega_{k}^{\prime}\)，论证类似。

